\label{sec:results}
\begin{figure}[t]
  \centering
  \includegraphics[width=1.2\columnwidth]{asset/pics/results/long_seq_blur.pdf}
  \caption{\textbf{Long Sequence Generation (24\,s Multi-Camera Rollout).} \modelname~produces temporally stable and view-consistent multi-camera videos over a 24\,s horizon, supporting streaming long-horizon simulation without catastrophic drift.}
  \label{fig:long_seq}
\end{figure}
This section evaluates \modelname~on its core capabilities as a controllable ego-centric multi-camera world model. We organize the results to highlight fine-grained controllability and generation quality, including: (i) ego action controllability, (ii) dynamic-agent and static-element controllability, (iii) long-horizon generation, (iv) multi-camera consistency, and (v) global appearance editing.

\subsection{Ego Action Controllability}
\label{subsec:action_ctrl}

Ego action controllability measures whether the generated future videos faithfully reflect the consequences of a given ego action sequence under a fixed initial scene context. As shown in Figure~\ref{fig:action_ctrl}, \modelname~produces photorealistic and physically plausible futures that closely follow the commanded actions. Conditioned on the same initial frame, changing the ego action sequence leads to correspondingly different, coherent rollouts, including smooth turns and lane changes. These results indicate that \modelname~internalizes a robust action-conditioned state transition in video space, enabling reactive and counterfactual rollouts for planning and policy evaluation.

\subsection{Dynamic and Static Element Controllability}
\label{subsec:ds_ctrl}

Beyond ego control, \modelname~supports fine-grained controllability over both \textbf{dynamic traffic agents} and \textbf{static road elements}, which is crucial for reproducible simulation and targeted stress testing. Figure~\ref{fig:ds_ctrl} visualizes this capability on \textbf{6\,s multi-camera} generations by overlaying projected conditions on the generated videos: {\color{green}green} boxes denote dynamic objects (\eg, vehicles), while {\color{red}red} and {\color{cyan}cyan} lines denote road boundaries and lane markings, respectively. The model preserves the specified agent locations/motions and anchors static layout cues consistently over time despite ego motion and viewpoint changes across different camera types (\eg, fisheye and narrow views). This demonstrates that dynamic and static controls can be enforced reliably and simultaneously, providing a controllable simulator interface for scenario editing and reproducible benchmarking.

\subsection{Long-Horizon Multi-Camera Generation}
\label{subsec:long_seq}

Figure~\ref{fig:long_seq} demonstrates \textbf{24\,s} long-horizon rollouts of \modelname~on multi-camera streams. The generated videos remain stable over extended horizons, maintaining coherent motion and appearance without catastrophic drift. Importantly, \modelname~can roll out beyond the short-clip setting used in Stage-I training, making it suitable for interactive simulation and closed-loop evaluation.

\subsection{Multi-Camera Consistency}
\label{subsec:mv_consistency}

Multi-camera world simulation requires both temporal coherence and cross-view alignment at each timestep. As illustrated in Figure~\ref{fig:long_seq} and Figure~\ref{fig:ds_ctrl}, \modelname~maintains strong \textbf{cross-view consistency}: dynamic objects exhibit consistent identity and motion across cameras, and static structures (\eg, lane topology and road boundaries) remain geometrically aligned across views. These qualitative results validate the effectiveness of our view--temporal modeling design for multi-camera generation.

\subsection{Global Appearance Editing via Text Prompts}
\label{subsec:style_transfer}

Finally, \modelname~supports controllable appearance editing through its text-conditioning branch, which primarily modulates global style while actions and scene elements are governed by their dedicated control branches. Figure~\ref{fig:style_transfer} shows C2V generations from the \textbf{front camera} where we keep ego actions and dynamic/static controls fixed, but vary the text prompt to edit global appearance, including locale (\eg, Germany vs.\ China), time of day (sunset vs.\ night), and weather (sunny vs.\ rainy). Notably, all frames are generated by \modelname, including the first frame, highlighting its ability to perform appearance-driven style transfer and controllable data synthesis.

% ---------------- Figures ----------------

\begin{figure}[t]
  \centering
  \includegraphics[width=1.0\columnwidth]{asset/pics/results/c2v_blur.pdf}
  \caption{\textbf{C2V for Scene Appearance Editing.} Keeping the ego actions and the dynamic/static elements fixed, we vary the text prompt to edit global appearance, including locale (\eg, Germany vs.\ China), time of day (sunset vs.\ night), and weather (sunny vs.\ rainy). All images are generated by \modelname, including the first frame.}
  \label{fig:style_transfer}
\end{figure}